\section{Positive path density near a saturated terminal}
\label{sec:positive-port-localization}

\begin{lemma}[A positive-density port lies near the far terminal]
\label{lem:positive-port-clustering}
Let $T$ be a normalized three-pole of rank $h$, with ports $a,b,c$.
Suppose $P_{ac}(T)=2^h$.
For every $0<\delta\le1$, the implication
\[
 P_{ab}(T)\ge\delta2^h
 \quad\Longrightarrow\quad
 \operatorname{dist}_T(b,c)
 \le2\log_2\!\left(\frac3{2\delta}\right)+3
\]
holds.
Moreover,
\begin{equation}\label{eq:one-full-product-tradeoff}
 P_{ab}(T)P_{bc}(T)\le7\,2^{3h/2}.
\end{equation}
In particular, $P_{ab}(T)\ge\delta2^h$ gives
$P_{bc}(T)\le(7/\delta)2^{h/2}$.
\end{lemma}

\begin{proof}
Suppress $b$ and use the saturated two-pole classification.
Orient its series chain from $a$ to $c$.
Suppose first that $b$ subdivides an edge inside a nonempty capped-ladder piece.
Let $p,q$ be its prefix and suffix position counts, and let $L,R$ be the sums of the ranks of earlier and later pieces.
The active piece has rank $p+q+d$, where $d\in\{1,2,3\}$ is the corresponding row parameter in the local subdivision table.
The prefix-cap estimates in Lemma~\ref{lem:one-full-tripole-small-adjacent} give
\[
 \frac{P_{ab}}{2^h}\le\frac32\,2^{-R-q/2},\qquad
 \frac{P_{bc}}{2^h}\le\frac32\,2^{-L-p/2}.
\]
The first inequality and $P_{ab}\ge\delta2^h$ imply
\[
 q+2R\le2\log_2\!\left(\frac3{2\delta}\right).
\]
Inside the active piece, a route from $b$ to its right terminal has length at most $q+3$.
Each later saturated piece has a terminal path whose length equals its rank.
The number of intervening bridges is at most the sum $R$ of those positive ranks.
Hence $\operatorname{dist}_T(b,c)\le q+2R+3$, proving the distance bound.

Multiplying the two normalized estimates gives
\[
 \frac{P_{ab}P_{bc}}{2^{2h}}
 \le\frac94\,2^{-(L+R)-(p+q)/2}
 \le\frac94\,2^{3/2}2^{-h/2}<7\,2^{-h/2}.
\]
Here $h=L+R+p+q+d$ and $d\le3$.

If the active piece is a digon, the normalized counts are $2^{-R}$ and $2^{-L}$, with $h=L+R+1$.
Its contribution to the distance from $b$ to the right terminal is at most two.
Both stated inequalities follow immediately.
If $b$ subdivides a bridge, the normalized counts are $2^{-R},2^{-L}$ with $h=L+R$.
The route from $b$ to $c$ has length at most $2R+1$.
Again both inequalities follow.
\end{proof}

\begin{theorem}[Three full arms and a positive fourth arm]
\label{thm:partial-five-star-return}
For each fixed $0<\delta\le1$, there is a constant $D_\delta$ with the following property.
Let $T$ be a normalized five-pole of rank $h$, with ports $c,a_1,a_2,a_3,a_4$.
Suppose
\[
 P_{ca_1}=P_{ca_2}=P_{ca_3}=2^h,
 \qquad P_{ca_4}\ge\delta2^h.
\]
Then the number of forests with odd-degree vertex set $\{a_1,a_2,a_3,a_4\}$ is at most
\[
 D_\delta2^{h/2}.
\]
\end{theorem}

\begin{proof}
First assume that $T$ is bridgeless.
Suppress $a_2,a_3$ and apply Lemma~\ref{lem:positive-port-clustering} to the ports $c,a_4,a_1$.
Restoring the two suppressed vertices increases the distance bound by at most two.
Thus
\[
 \operatorname{dist}_T(a_4,a_1)
 \le2\log_2\!\left(\frac3{2\delta}\right)+5=:d_\delta.
\]
Full terminal-path saturation forces every $c$--$a_i$ path, for $i\le3$, to have at least $h$ edges.
The triangle inequality therefore gives
\[
 \operatorname{dist}_T(c,a_i)\ge h-d_\delta
 \qquad(1\le i\le4).
\]

Use the bridgeless proof of Theorem~\ref{thm:far-terminal-forests}.
That proof uses the high path densities only to obtain a far-tail distance bound and a bounded residual rank for one normalized two-pole.
The displayed distance bounds supply the first condition, with constant $d_\delta$.
The full pair $c,a_1$ supplies the second condition, with additive-stability constant $C_{1/2}$.
Choose the layer cutoff using the larger of these two constants.
All marked sets still have bounded size depending only on $\delta$.
The same single-minimum-layer encoding proves the claimed bound.

Now suppose that $T$ has a bridge.
Both shores contain cycles, by the degree count used in Theorem~\ref{thm:far-terminal-forests}.
The three full pairs force $c$ onto one shore and $a_1,a_2,a_3$ onto the other.
Neither $c$ nor those three leaf ports can be a bridge endpoint.
Call the two ranks $h_L,h_R$, with $c$ on the left.

If $a_4$ is a bridge endpoint, every $c$--$a_4$ path avoids the cyclic part of the right shore.
Thus $P_{ca_4}\le2^{h_L}$ and $2^{h_R}\le1/\delta$.
If $a_4$ is the left endpoint, each counted forest uses the bridge, has no left-shore edge, and gives a four-terminal forest on the right.
If $a_4$ is the right endpoint, each counted forest stays on the right.
In both cases, the forest count is at most $2^{h_R}\le1/\delta$.

Otherwise, if $a_4$ lies on the right, retain the right shore and make its bridge attachment the new center.
Full path counts factor across the bridge, so all three full pairs remain full there.
The fourth normalized density remains at least $\delta$.
Every counted forest stays in that shore.
Induction on rank proves the result.

Finally, suppose that $a_4$ lies on the left and is not the bridge endpoint.
The left shore is a normalized three-pole with ports $c,a_4,z$, where $z$ is its bridge attachment.
Full path factorization gives $P_{cz}=2^{h_L}$.
Every counted forest uses the bridge exactly once.
Its left part is an $a_4$--$z$ path, and its right part is a four-terminal forest.
The latter has at most $2^{h_R}$ choices.
Equation~\eqref{eq:one-full-product-tradeoff} gives
\[
 P_{a_4z}\,P_{ca_4}\le7\,2^{3h_L/2}.
\]
Since $P_{ca_4}\ge\delta2^{h_L+h_R}$, the total forest count is at most
\[
 2^{h_R}P_{a_4z}
 \le\frac7\delta\,2^{h_L/2}
 \le\frac7\delta\,2^{h/2}.
\]
Increasing $D_\delta$ covers every bridge case.
\end{proof}

This theorem excludes a five-pole profile with three full star arms, a fixed positive fourth arm, and a full four-leaf forest sector.
It does not impose a fixed lower bound on the fourth arm when that arm tends to zero with rank.
